% Generated mechanically from frozen input p01/ko/introduction.tex.
% Do not edit this generated file; edit the bound locale source instead.

% OK, start here.
%

\setcounter{chapter}{0}
\chapter{서론}



\phantomsection
\label{introduction-section-phantom}

\begin{verbatim}
Copyright (C) 2005 -- 2025 Johan de Jong
Permission is granted to copy, distribute and/or modify this
document under the terms of the GNU Free Documentation License,
Version 1.2 or any later version published by the Free Software
Foundation; with no Invariant Sections, no Front-Cover Texts,
and no Back-Cover Texts. A copy of the license is included in
the section entitled "GNU Free Documentation License".
\end{verbatim}



\section{개요}
\label{introduction-section-overview}

\noindent
Laumon과 Moret--Bailly의 저서(\cite{LM-B} 참조), 그리고 Fulton 등이 집필 중인
저작과 더불어, 대수 스택과 이를 정의하는 데 필요한 대수기하학을 다루는 오픈 소스
교재가 필요하다고 우리는 생각한다. Stacks Project는 가환대수학에서 출발하여
스킴 이론과 대수공간 이론을 거쳐 대수 스택 이론의 포괄적인 기초를 세움으로써
이 역할을 수행하고자 한다.

\medskip\noindent
이 자료는 온라인으로 읽히는 것을 전제로 한다. 여러 페이지에 흩어진 내부 참조로
빠르게 이동할 수 있는 하이퍼링크가 핵심 기능이기 때문이다. 브라우저에 내장된
PDF 또는 DVI 뷰어를 사용하면 파일 사이의 링크도 작동할 것이다.

\medskip\noindent
이 프로젝트는 공동 작업이며, 독자의 참여를 권한다. 읽는 중에 발견한 오자나 오류,
또는 제안, 추가 자료, 예가 있다면
\href{mailto:stacks.project@gmail.com}{stacks.project@gmail.com}으로
이메일을 보내 주기 바란다. 모든 원본 파일이 들어 있는 tarball을 내려받아 압축을 풀고
make를 실행하면 DVI 또는 PDF 뷰어로 로컬에서 볼 수 있다. LaTeX 파일을 자유롭게
편집하고 개선 사항을 이메일로 보내도 좋다.


\section{기여 내역}
\label{introduction-section-attribution}

\noindent
이 작업의 범위는 매우 넓어서, 여기서 설명하려는 이론의 발전에 기여한 모든 수학자에게
정확하고 간결하게 공로를 돌리는 일은 벅찬 과제이다. 언젠가 공동체의 관심이 충분히
커져서, 각 장마다 역사적 주석을 담은 절을 써 줄 기여자를 찾을 수 있기를 바란다.

\medskip\noindent
이 작업에 기여한 사람들은 책 판본의 표제지와
\href{https://stacks.math.columbia.edu/tex/CONTRIBUTORS}{온라인 명단}에
나열되어 있다. 여기서는 주요 기여의 일부를 소개한다.
\begin{enumerate}
\item Jarod Alper는 대수 스택 관련 문헌을 논하는 장을 기여하였다.
『문헌 안내』 제 \ref{guide-section-short-introductions} 절을 참조하라.
\item Bhargav Bhatt는 스킴의 에탈 사상에 관한 장의 초고를 작성하였다.
『에탈 사상』 제 \ref{etale-section-introduction} 절을 참조하라.
\item Bhargav Bhatt는
『대수학 심화』 제 \ref{more-algebra-section-formal-glueing} 절의 초고를 작성하였다.
\item Kiran Kedlaya는
『내림』 제 \ref{descent-section-descent-universally-injective} 절의 초고를 기여하였다.
\item 다음 부분의 초고는
\begin{enumerate}
\item 『대수학』 제 \ref{algebra-section-oka-families} 절,
\item 『단사 대상』 제 \ref{injectives-section-baer} 절, 그리고
\item 체에 관한 장, 즉
『체』 제 \ref{fields-section-introduction} 절이다.
\end{enumerate}
이들은 Akhil Mathew 등의 호의로
\href{https://math.uchicago.edu/~amathew/cr.html}{The CRing Project}에서
가져왔다.
\item Alex Perry는 사영 가군과 Mittag--Leffler 가군에 관한 내용을 작성하였으며,
『대수학』 정리 \ref{algebra-theorem-ffdescent-projectivity}의 증명도 포함된다.
\item Alex Perry는 Schlessinger와 Rim의 방식에 따른 변형 이론의 장을 작성하였다.
『형식적 변형 이론』 제 \ref{formal-defos-section-introduction} 절을 참조하라.
\item Thibaut Pugin, Zachary Maddock, Min Lee는 에탈 코호몰로지의 장과
트레이스 공식의 장의 바탕이 된 강의 노트를 작성하였다. 『에탈 코호몰로지』 제
\ref{etale-cohomology-section-introduction} 절과
『트레이스 공식』 제 \ref{trace-section-introduction} 절을 참조하라.
\item David Rydh는 많은 유익한 의견을 제공하고 여러 오류를 지적했으며,
잔여 gerbe에 관한 내용에 핵심적으로 기여하였다. 또한
『공간에서의 군자 심화』 제
\ref{spaces-more-groupoids-section-finite} 절과
\ref{spaces-more-groupoids-section-etale-localize} 절의 내용을 처음 제안하였다.
\item Burt Totaro는 『예』 제
\ref{examples-section-non-descending-property-projective} 절,
\ref{examples-section-non-effective-descent-projective} 절과
『스택의 성질』 제 \ref{stacks-properties-section-dimension} 절을 기여하였다.
\item pro-에탈 코호몰로지의 장, 즉
『pro-에탈 코호몰로지』 제 \ref{proetale-section-introduction} 절은
Bhargav Bhatt와 Peter Scholze의 논문에서 가져왔다.
\item Bhargav Bhatt는 『예』 제
\ref{examples-section-non-algebraic-hom-stack} 절과
\ref{examples-section-flat-not-colimit-flat-finitely-presented} 절을 기여하였다.
\item Ofer Gabber는 오류를 발견하고 수정 사항을 기여했으며, 다음 내용도 기여하였다.
『대수다양체』 보조정리 \ref{varieties-lemma-image-connected-component},
『형식적 공간』 보조정리 \ref{formal-spaces-lemma-completion-countably-indexed},
『군자 심화』 제 \ref{more-groupoids-section-ind-quasi-affine} 절의 내용,
『공간의 성질』 제 \ref{spaces-properties-section-fpqc} 절의 주된 결과, 그리고
『평탄성 심화』 명제
\ref{flat-proposition-finite-type-injective-into-flat-mod-m}의 증명이다.
\item J\'anos Koll\'ar는
『대수학』 보조정리 \ref{algebra-lemma-hart-serre-loc-thm}와
『국소 코호몰로지』 명제 \ref{local-cohomology-proposition-kollar}를 기여하였다.
\item Kiran Kedlaya는
『대수학 심화』 제 \ref{more-algebra-section-beauville-laszlo} 절의 초고를 작성하였다.
\item Matthew Emerton, Toby Gee, Brandon Levin은 두꺼워짐에 관한 몇 가지 결과,
특히 『스택 사상 심화』의 보조정리
\ref{stacks-more-morphisms-lemma-reduced-diagonal},
\ref{stacks-more-morphisms-lemma-thickening-diagonals}, 그리고
\ref{stacks-more-morphisms-lemma-thickening-properties}를 기여하였다.
\item Lena Min Ji는
『대수학 심화』 제 \ref{more-algebra-section-principal-radical-ideals} 절의 초고를 작성하였다.
\item Matthew Emerton과 Toby Gee는 『스택의 기하학』 제
\ref{stacks-geometry-section-multiplicities} 절과
\ref{stacks-geometry-section-dimension-of-algebraic-stacks} 절의 초고를 작성하였다.
\end{enumerate}
